\DocumentMetadata{
  pdfversion=2.0,
  pdfstandard=ua-2,
  lang=en-US,
  tagging=on,
  tagging-setup={math/setup=mathml-SE,tabsorder=structure}
}
\documentclass[11pt,letterpaper]{article}
\usepackage[margin=0.68in,includefoot]{geometry}
\usepackage{newcomputermodern}
\usepackage{amsmath,amscd,mathtools}
\providecommand{\square}{\mdlgwhtsquare}
\usepackage[hidelinks]{hyperref}
\hypersetup{
  pdftitle={Complex Analysis PhD Exam, May 2021},
  pdfauthor={University of Florida Department of Mathematics},
  pdfsubject={Graduate mathematics examination},
  pdfdisplaydoctitle=true,
  bookmarksopen=true
}
\setcounter{secnumdepth}{0}
\setlength{\parindent}{0pt}
\setlength{\parskip}{0.28em}
\setlength{\emergencystretch}{3em}
\setlength{\leftmargini}{2.15em}
\setlength{\leftmarginii}{2.65em}
\setlength{\leftmarginiii}{3em}
\renewcommand{\labelenumi}{\textbf{\theenumi.}}
\pagestyle{empty}
\raggedbottom
\allowdisplaybreaks
\newenvironment{examproblems}[1]{%
  \begin{enumerate}%
  \setcounter{enumi}{#1}%
  \setlength{\topsep}{0.35em}%
  \setlength{\partopsep}{0pt}%
  \setlength{\itemsep}{0.48em}%
  \setlength{\parsep}{0.18em}%
}{\end{enumerate}}
\newenvironment{examsubparts}{%
  \begin{enumerate}%
  \setlength{\topsep}{0.18em}%
  \setlength{\partopsep}{0pt}%
  \setlength{\itemsep}{0.16em}%
  \setlength{\parsep}{0.1em}%
}{\end{enumerate}}
\newenvironment{examinstructions}{%
  \par\begingroup\itshape
}{%
  \par\endgroup\vspace{0.18em}
}
\makeatletter
\renewcommand\section{\@startsection{section}{1}{0pt}%
  {-1.25ex plus -.25ex minus -.1ex}%
  {0.55ex plus .1ex}%
  {\normalfont\large\bfseries}}
\renewcommand{\@maketitle}{%
  \newpage\null\vskip -1.2em%
  {\footnotesize\bfseries
    \href{https://www.ufl.edu/}{University of Florida}, \href{https://math.ufl.edu/}{Department of Mathematics}\par}%
  \vskip 0.18em%
  \tagstructbegin{tag=Title}%
    {\LARGE\bfseries\href{https://gma.math.ufl.edu/exams/complex-analysis-phd/}{Complex Analysis PhD Exam}, May 2021\par}%
  \tagstructend%
  \vskip 0.35em%
  \tagmcbegin{artifact}%
    \hrule height 1pt%
  \tagmcend%
  \vskip 0.55em%
}
\makeatother
\title{Complex Analysis PhD Exam, May 2021}
\author{}
\date{}
\begin{document}
\maketitle
\thispagestyle{empty}
\begin{examinstructions}
Solve five of the following problems. Please explain your work in a careful and logical way. You can use any theorem, identity, functional equation proved in class during Fall 2020 and Spring 2021. Please make sure to state precisely any result or theorem you are using in your solutions. Failing to do so may result in a lower grade.
\end{examinstructions}
\begin{examproblems}{0}
\item
Let \(\Gamma(z)\) be the Euler gamma function. Recall this is a meromorphic function over \(\mathbb C\) we defined via the infinite product 
\[
\Gamma(z)=\frac{e^{-\gamma z}}{z}\prod_{k=1}^{\infty}\left(1+\frac{z}{k}\right)^{-1}e^{z/k},
\]
 where~\(\gamma\) is the Euler constant. Compute the integral 
\[
\int_{C_{5/2}(0)}\Gamma(z)\,dz,
\]
 where \(C_{5/2}(0)\) is the circle of radius \(5/2\) centered at the origin oriented counterclockwise.
\item
Show that the polynomial \(z^7-4z^3+z-1\) has three zeros inside the unit circle (counted with multiplicities).
\item
Let~\(U\) be a simply connected open subset of \(\mathbb C\) such that \(U\ne\mathbb C\). Show that the group of analytic automorphisms of~\(U\), say \(\operatorname{Aut}(U)\), acts transitively on~\(U\). What happens if \(U=\mathbb C\)?
\item
Find all entire functions \(f:\mathbb C\to\mathbb C\) such that 
\[
f(f(z))=z,
\]
 for all \(z\in\mathbb C\).
\item
Show that if \(f:\mathbb C\to\mathbb C\) is a continuous function such that~\(f\) is holomorphic on \(\mathbb C\setminus[-1,1]\) (the open set obtained by removing the closed interval \([-1,1]\) from the complex plane), then~\(f\) is necessarily analytic everywhere.
\item
Let \(P(z)\) be a polynomial with complex coefficients. Show that if the degree of \(P(z)\) is at least two then the sum of residues of \(1/P(z)\) (over all zeros of \(P(z)\)) is zero. What happens if \(P(z)\) has degree one?
\item
Consider the open subset of \(\mathbb C\) 
\[
E:=\{z\in\mathbb C\mid |z|>1\}.
\]
 Does there exist an analytic isomorphism between~\(E\) and \(\mathbb C^*:=\mathbb C\setminus\{0\}\)?
\item
Let \(P_n(z)\) be a polynomial of order~\(n\) such that 
\[
|P_n(z)|\le M,\qquad M>0,
\]
 for \(|z|=1\). Show that 
\[
|P_n(z)|\le M|z|^n
\]
 for \(|z|\ge 1\).
\item
Show that there is no meromorphic function over \(\mathbb C\) such that 
\[
f(z+1)=f(z),\qquad f(z+i)=f(z)
\]
 for all \(z\in\mathbb C\), whose only poles are simple poles at the points 
\[
m+ni,\qquad m,n\in\mathbb Z.
\]
 What if those points are assumed to be poles of order two and~\(f\) has no other poles?
\end{examproblems}
\end{document}
